% Generated draft: requires mathematical and source review.
\documentclass[11pt]{article}
\usepackage[margin=25mm]{geometry}
\usepackage{fontspec}
\setmainfont{Latin Modern Roman}
\usepackage{amsmath,amssymb,mathtools,unicode-math}
\setmathfont{Latin Modern Math}
\usepackage{enumitem,needspace,xurl,hyperref}
\setlength{\parindent}{0pt}
\setlength{\parskip}{5pt}
\setlength{\emergencystretch}{3em}
\newcommand{\N}{\mathbb{N}}
\newcommand{\Z}{\mathbb{Z}}
\newcommand{\Q}{\mathbb{Q}}
\newcommand{\R}{\mathbb{R}}
\newcommand{\C}{\mathbb{C}}
\newcommand{\F}{\mathbb{F}}
\newcommand{\A}{\mathbb{A}}
\newcommand{\PP}{\mathbb{P}}
\newcommand{\E}{\mathbb{E}}
\newcommand{\eps}{\varepsilon}
\newcommand{\dd}{\,\mathrm{d}}
\begin{document}
\section*{UnsolvedMath ID: 30000477}
\textbf{Source problem code:} catalogue entry
\subsection{Definitions and mathematical statement}
Fix the source parameter $\lambda$, and suppose $M$ is an $AD_{\R}$-like model at $\lambda$. Put $\theta=\Theta_M$ and define the fixed-parameter, fixed-$\Theta$ collection
\[
 \mathcal C_{\lambda,\theta}
 =\{N:N\text{ is }AD_{\R}\text{-like at }\lambda,
                          \ \Theta_N=\theta\},\qquad
 M_0=\bigcap_{N\in\mathcal C_{\lambda,\theta}}N.
\]
The family is nonempty since it contains $M$. All models and the intersection must be interpreted inside the same ambient framework from the report. If this is a class-sized family, the displayed intersection is the corresponding common-membership class; the source's set/class and height conventions are not supplied.

The Minimum Model Conjecture at $\lambda$ asks whether
\[
 M_0\text{ is }AD_{\R}\text{-like at }\lambda
 \qquad\text{and}\qquad \Theta_{M_0}=\theta.
\]
Both conclusions are required. It is not enough that the intersection is contained in every candidate, and intersecting models with different $\Theta$ values is not the source assertion. If the conclusion holds, $M_0$ is the minimum member under actual inclusion of this fixed collection, not merely a model unique up to isomorphism.

In the usual determinacy-model convention, $\Theta_N$ is the supremum of ordinals that are surjective images of the reals of $N$ by maps belonging to $N$:
\[
 \Theta_N=\sup\{\alpha:\exists f\in N,
                         \ f:\R^N\twoheadrightarrow\alpha\}.
\]
The exact internal/ambient interpretation is part of the omitted framework. The name $AD_{\R}$ ordinarily concerns determinacy of length-$\omega$ games with real moves, but \emph{$AD_{\R}$-like at $\lambda$} is a stronger or different source-defined predicate, not simply the instruction to replace it by ``satisfies $AD_{\R}$''. Its model axioms, shared-real/ordinal conditions and relation to $\lambda$ are absent and are not invented.

The supplied literature triage checked 21 August 2026 calls this conjectural and says it depends on an unresolved weak-square-root program. No definition or exact logical reduction for that program is supplied. The dated claim is preserved without asserting an independently verified current status or proving that intersections retain the required model properties.
\subsection{Short English statement}
For a fixed $\lambda$ and fixed $\Theta$, does the intersection of all source-defined $AD_{\R}$-like models remain such a model and retain exactly that $\Theta$? The technical meaning of like-at-$\lambda$ and the ambient model framework are missing. Both closure under intersection and preservation of $\Theta$ are required.
\subsection{Sources}
\begin{itemize}
\item Local UnsolvedMath registry, permanent ID 30000477, exact source code catalogue entry; Minimum $AD_{\R}$-Like Model at $\lambda$. Proposed year 2006, status open. The corrected\_verified statement restores the complete fixed-$\Theta$ intersection and both conclusions; claimed report inspection is source metadata.
\item Supplied report citation: \emph{Feinstrukturtheorie und Innere Modelle (2007)}, Oberwolfach Reports, DOI 10.4171/owr/2006/20, pp. 1215--1244, \url{https://doi.org/10.4171/owr/2006/20}. Title year differs from DOI/proposed year.
\item Supplied literature triage checked 2026-08-21: ``original OWR report 20/2006'', \url{https://doi.org/10.4171/owr/2006/20}; ``OWR text reproducing the open-problem statement'', \url{https://ems.press/content/serial-article-files/46048}. It reports the Minimum Model Conjecture and its weak-square-root-program dependence. This dated OWR-derived curation requests independent verification of time-sensitive status and cited claims.
\end{itemize}
\end{document}
